\chapter{Badania własne i analiza wyników}
\label{ch:badania}

% Tabele i wykresy generuje: python -m anomaly.experiments (skopiuj results/tables -> tables/, results/figures -> figures/)
% \input{tables/tab_podzial.tex}   \input{tables/tab_wyniki_glowne.tex}   \input{tables/tab_przedzialy.tex}
% \input{tables/tab_niu.tex}   \input{tables/tab_niezbalansowanie.tex}   \input{tables/tab_okna.tex}
\section{Metodologia eksperymentu}
\todo{Podział danych, walidacja, hiperparametry, powtórzenia (seedy), sprzęt.}
\section{Skuteczność algorytmów uczenia nadzorowanego}
\todo{Tabela metryk, macierze pomyłek.}
\section{Skuteczność algorytmów uczenia nienadzorowanego}
\todo{Tabela metryk.}
\section{Wpływ sposobu radzenia sobie z niezbalansowaniem klas}
\todo{Wagi klas vs SMOTE.}
\section{Porównanie wyników}
\todo{Krzywe ROC/PR, wpływ progu, porównanie z \cite{niu2019}.}
\section{Analiza statystyczna wyników}
\todo{Powtórzenia i bootstrap; przedziały ufności.}
\section{Ocena wydajności systemu i ograniczenia}
\todo{Opóźnienie oceny, ograniczenia zbioru (cechy PCA, dwa dni, brak dryfu).}
